%###########################################################
%#
%#
%#
%###########################################################
\section{Research Project}

Fuel cells are electrochemical energy conversion devices that directly transform chemical energy into electrical energy without the need for combustion processes, making them essential technologies for the global energy transition toward sustainable power generation systems \cite{Nigel2017}. Unlike conventional thermal devices, fuel cells rely on electrochemical reactions occurring at electrodes separated by an ion-conducting electrolyte, enabling high efficiency and low environmental impact \cite{Gurbinder2016}.

Fuel cells convert chemical energy directly into electrical energy through a redox reaction. Like thermal devices, they operate continuously as long as fuel and oxidant are supplied. Because of this nature, fuel cells are capable of continuously producing electricity with high conversion efficiency.

Several types of fuel cell technologies exist, each defined primarily by the electrolyte material and operating temperature range \cite{Gurbinder2016,skinnerd:2008}. Examples include polymer electrolyte membrane fuel cells, alkaline fuel cells, and solid oxide fuel cells. Despite differences in materials and temperature regimes, all fuel cell technologies share common characteristics such as high energy conversion efficiency, absence of moving mechanical parts, modular scalability, and low or near-zero pollutant emissions during operation.

Applications of fuel cells extend to a wide technological spectrum, from portable electronic devices and backup power systems to transportation and large-scale distributed electricity generation \cite{BinZhu2020}. These technologies contribute to reducing dependence on fossil fuels, lowering greenhouse gas emissions, improving energy security, and supporting the development of a hydrogen-based energy infrastructure.

Growing global concern regarding energy security, atmospheric pollution and greenhouse gas emissions has intensified research into environmentally sustainable energy technologies. Among the various fuel cell systems, Solid Oxide Fuel Cells (SOFC) have attracted considerable scientific and technological interest due to their high efficiency, fuel flexibility, and long-term stability \cite{Kevin2003}. SOFCs can operate using hydrogen, carbon monoxide, methane, and other hydrocarbon-based fuels, reducing the need for highly purified hydrogen and enabling integration with existing fuel infrastructures.

The project focuses on  symmetric Intermediate-Temperature Solid Oxide Fuel Cells \cite{JuanCarlos2011,KejunZhu2022,Muhammad2022,Liliana2016}, which operate typically between 500 and 700°C. Operating in this temperature range reduces thermal degradation, expands material choices, and lowers system costs while maintaining sufficient ionic conductivity and catalytic activity. Because of their symmetric configuration and reversible electrochemical behavior, these devices can also function as Solid Oxide Electrolyzer Cells (SOEC), enabling efficient electrolysis of water \cite{Meng2008}. When powered by renewable electricity, SOEC systems can produce green hydrogen, a key energy carrier in decarbonized energy systems.

As shown in Figure \ref{sofc}, a Solid Oxide Fuel Cell device is composed of a cathode, an anode, and a dense ceramic electrolyte. At the cathode, oxygen reduction reactions (ORR) convert molecular oxygen into oxide ions. Perovskite oxides such as lanthanum strontium manganite exhibit excellent electrocatalytic performance for this reaction \cite{Sacanell2017,Diego2017,JACOBSON2012478}. The anode promotes fuel oxidation (or Hydrogen Oxidation Reaction, HOR), commonly employing nickel yttria stabilized zirconia cermets capable of hydrogen oxidation and internal reforming of methane-containing fuels \cite{Nigel2017,skinnerd:2008}. The electrolyte is typically a dense yttria-stabilized zirconia ceramics, which provides high oxide-ion conductivity and chemical stability at elevated temperatures \cite{Nigel2017,skinnerd:2008}.

\begin{figure}[h]
\centering
    \begin{overpic}[width=.75\linewidth,keepaspectratio,angle=0,trim=0.cm 0.cm 0.cm 0.cm, clip]{Figures/SOFC.png}
    \put(45,85){\Large \bf \textcolor{Blue}{Cathode}}
    \put(62,60){\Large \bf \textcolor{Black!90}{Electrolyte}}
    \put(-3,55){\Large \bf \textcolor{Orange}{Anode}}

    \put(30,75){\Large \bf \textcolor{White}{1/2O$_2$+2e$^-$ $\to$ 2O$^{2-}$ }}
    
    \put(41,63,55){\Large \bf \textcolor{Black}{ $\downarrow$ }}
    \put(45,62.5){\Large \bf \textcolor{Black}{ O$^{2-}$ }}
    
    \put(30,35){\Large \bf \textcolor{White}{H$_2$+O$^{2-}$ $\to$ H$_2$O+2e$^-$}}
    
    \end{overpic}
    
    \caption{SOFC device scheme showing the Cathode Oxygen Reduction Reaction (ORR), and the Anode Hydrogen Oxidation Reaction (HOR). The ionic conductivity ($\sigma_{ion}$) of the electrolyte is governed by the migration of oxygen vacancies ($V_O ^{\cdot\cdot}$) in the crystal lattice, which follows the Arrhenius law dependence on temperature \cite{Nigel2017}. Figure addapted from \cite{HU2023}. }
    
    \label{sofc}
\end{figure}


A major limitation of conventional SOFC technology is that peak efficiency occurs at temperatures near 900–1000$^\circ$C \cite{Nigel2017,Gurbinder2016,skinnerd:2008}. Such high temperatures restrict material selection for interconnects and structural components and accelerate degradation mechanisms such as thermal expansion mismatch, phase instability, and chemical inter diffusion. And materials such as lanthanum, strontium, and chromium are often required, which are costly. These challenges have motivated extensive research into materials capable of efficient operation at intermediate temperatures.

Despite significant advances in intermediate-temperature SOFC materials, several \linebreak state-of-the-art cathode materials rely on elements such as lanthanum, cobalt, and cerium, which present economic, and environmental concerns. This project proposes the exploration of alternative mixed ionic-electronic conductors based on more abundant and lower-cost elements such as iron and molybdenum. Candidate materials include perovskite-structured compounds such as (La,Sr)(Fe,Mo)O$_{3-\delta}$, (Ba,Sr)(Fe,Mo)O$_{3-\delta}$, (La,Sm)(Fe)O$_{3-\delta}$, and  (La,Sm)(Fe,Mo)O$_{3-\delta}$. These materials have shown promising electrochemical behavior in related studies \cite{GIL2024,MEJIAGOMEZ2020} and may provide competitive catalytic performance while reducing cost and environmental impact.

\hl{Would we like to expand the set of candidate materials for SOFC applications and provide technical justification for their inclusion?}\dotfill

\dotfill

\dotfill

\dotfill

\dotfill

In parallel, Solid Oxide Electrolyzer Cell (SOEC) systems will be developed using cobalt-free oxygen electrode materials based on (Ln,Sr)Fe$_{0.8}$Cu$_{0.2}$O$_{3-\delta}$ perovskite oxides, where Ln represents lanthanide elements such as La, Nd, Pr, or Gd. Among these, LaSFCu compositions have demonstrated excellent reversible performance in SOFC/SOEC operation. Composite materials incorporating ceria-based fluorites, such as LSFCu-CeO$_2$ nanocomposites with ceria contents, will also be investigated due to their enhanced oxygen transport and catalytic activity \cite{Natasha2026}.

\hl{Would we like to expand the set of candidate materials for SOEC applications and provide technical justification for their inclusion?}\dotfill

\dotfill

\dotfill

\dotfill

\dotfill

Another topic proposed in this project is the valuation and development of nano catalysts materials for the steam reforming of ethanol (SRE). Bio-ethanol energy is a renewable energy resource, and its utilization can reduce pollution \cite{FENG2023}.  Bio-ethanol can be obtained from the fermentation of various raw materials containing sucrose, and cellulosic biomass \cite{Tripodi2018}. It also has high hydrogen content, non-toxicity, easy storage, and processing. Therefore, producing hydrogen from ethanol reforming process has gained widespread attention. The ethanol reforming process can be achieved by different oxidants \cite{FENG2023} including H$_2$O for the process called steam reforming of ethanol. 

Commonly used SRE catalysts are supported catalysts with active group VIII metals and Cu, supported by thermostable oxide ceramics such as CeO$_2$, Al$_2$O$_3$, SiO$_2$ ZrO$_2$, etc. Considering the results from previous publications by the group \cite{}, we will work on the design of novel nano-structured Ni-Cu bimetallic materials as catalysts for steam reforming of ethanol and will adjust their composition and morphology to optimize reactivity and selectivity for this reaction.

\hl{Would we like to expand the set of candidate materials for SRE applications and provide technical justification for their inclusion?}\dotfill

\dotfill

\dotfill

\dotfill

\dotfill

The synthesized materials will undergo comprehensive characterization to determine structural, morphological, and redox properties. Techniques will include X-ray diffraction, and electron microscopy. Catalytic and electrocatalytic behavior will be evaluated using gas chromatography under controlled reaction conditions. These laboratory measurements will be performed at the CNPEM support facilities and will be complemented by in-situ and operando synchrotron based X-ray absorption spectroscopy experiments to directly observe structural and electronic changes during operation.

The research group possesses strong expertise in nanomaterial synthesis and preparation using multiple processing routes, including sol-gel synthesis, and combustion methods \cite{Diego2011,Lucia2016,Diego2024}. Those expertise ensures the feasibility of producing high-quality materials and overcoming synthesis challenges associated with complex perovskite systems and nanocomposite electrodes.

A central component of this project is the development of specialized experimental instrumentation compatible with synchrotron characterization. The project will leverage the capabilities of the QUATI beamline for nanoscale investigation of functional materials \cite{Santiago2023}. The methodology involves experimental engineering cells capable of operating under realistic catalytic and electrochemical conditions while allowing X-ray transmission and precise environmental control. These measurement cells will enable direct observation of catalytic reactions, electrocatalytic processes, and ionic transport phenomena in nanostructured materials during operation, establishing a powerful experimental framework for understanding and optimizing next-generation solid oxide electrochemical technologies.

\hl{Include also Time Differential Perturbed Angular Correlation technique for material characterization}\dotfill

\dotfill

\dotfill

\dotfill

\dotfill

Finally, the integration of advanced materials development techniques, catalytic hydrogen production via steam reforming of ethanol, and operando synchrotron characterization represents a comprehensive approach to addressing key scientific and technological challenges in electrochemical energy conversion systems, and will be the focus of this proposal.
